\subsection{纤维积与逆像}

5.11.1. 设 F 是一个 Banach 空间，\((\mathrm{E}_{i})_{i\in I}\) 是有限个 Banach 空间组成的族，\(\mathbf{f}=(f_{i})_{i\in I}\) 是由从 \(f_{i}\) 到 F 的连续线性映射 \(E_{i}\) 组成的族。令 E 为各 \(E_{i}\) 的乘积，令 f 为从 E 到 \(F^{I}\) 的连续线性映射，即各 \(f_{i}\) 的乘积。最后，令 D 为 \(F^{I}\) 的闭子空间，它由满足 \((y_{i})_{i\in I}\) 与 i 无关的 \(y_{i}\) 构成（\(F^{I}\) 的“对角线”）。称族 f 是横截的，如果将 f 与从 \(F^{I}\) 到商空间 \(F^{I}/D\) 的典范投影复合所得的连续线性映射是满射，且其核 \(f^{-1}(D)\) 具有拓扑补空间。

若空间 \(E_{i}\) 和 F 都是有限维的，则族 f 是

横截的，当且仅当有：

\[
\operatorname{codim} \left(\bigcap_ {i} \operatorname{Im} f _ {i}\right) = \sum_ {i} \operatorname{codim} \left(\operatorname{Im} f _ {i}\right)
\]

若 \(I = \{1, 2\}\)，则有序对 \((f_1, f_2)\) 横截，当且仅当映射

\[
f _ {1} + f _ {2}: \mathbf {E} _ {1} \oplus \mathbf {E} _ {2} \rightarrow \mathbf {F}
\]

是满射且其核具有拓扑补空间（如果 \(E_{1}\) 和 \(E_{2}\) 是有限维的，这等价于说

\[
\operatorname{Im} f _ {1} + \operatorname{Im} f _ {2} = \mathrm{F}).
\]

5.11.2. 设 S 是一个流形，\((X_{i})_{i\in I}\) 是有限个流形组成的族，\(\mathbf{f}=(f_{i})_{i\in I}\) 是由从 \(f_{i}\) 到 S 的态射 \(X_{i}\) 组成的族。令 P 为各 \(X_{i}\) 的乘积 X 的子集，它由满足 \((x_{i})_{i\in I}\) 与 i 无关的点 \(f_{i}(x_{i})\) 构成。令 \(x\in P\)，并令 \(y=f_{i}(x_{i})\in S\)。若映射 \(T_{x}(f_{i})\) 构成取值于 Banach 空间 \(T_{y}(S)\) 的连续线性映射的横截族，则称族 f 在 P 的点 x 处横截。若 f 在 P 的每一点处都横截，则称 f 是横截的。

若 f 是横截的，则 P 是 X 的子流形，称为在 S 上方（相对于族 f）的各 \(X_{i}\) 的纤维积，记为 \(\prod_{i\in I}X_{i}\)（例如当 \(X_{1}\times_{S}X_{2}\) 时也简写为 \(I=\{1,2\}\)）。对于 P 的每一点 \(x=(x_{i})\)，切空间 \(T_{x}(\prod_{S}X_{i})\) 是 \(\prod T_{x}(X_{i})\) 的子空间，由满足 \(t=(t_{i})\) 与指标 i 无关的切向量 \(T_{x_{i}}(f_{i}).t_{i}\) 组成。

若 \(f_{1}: X_{1} \to Y\) 是浸没，\(f_{2}: X_{2} \to Y\) 是态射，则有序对 \((f_{1}, f_{2})\) 是横截的。

5.11.3.（“纤维积的泛性质”）设 \(\mathbf{f} = (f_i)_{i \in I}\) 是态射 \(f_i: X_i \to S\) 的横截族，P 是在 S 上各 \(X_i\) 的纤维积；对于每个 \(i \in I\)，以 \(\pi_i\) 表示从 P 到 \(X_i\) 的态射，它由将 X 到 \(X_i\) 的投影限制在 P 上得到；于是 \(f_i \circ \pi_i\) 是从 P 到 S 的态射，且与 \(i\) 无关。设 T 是一个流形，\(g_i: T \to X_i\) 是流形的态射，并且 \(f_i \circ g_i\) 是从 T 到 S 的态射，且与 \(i \in I\) 无关；则存在唯一的从 T 到 P 的态射 \(h\)，使得对于所有 \(g_i = \pi_i \circ h\)，都有 \(i \in I\)。

5.11.4.（“纤维积的结合律”）设 \(\mathbf{f} = (f_i)_{i \in I}\) 是流形态射 \(f_i: \mathbf{X}_i \to \mathbf{S}\) 的有限族，\((\mathrm{J}_{\lambda})_{\lambda \in \Lambda}\) 是 I 的一个划分。假设对 \(\lambda\) 中每个 \(\Lambda\)，族 \(\mathbf{f}_{\lambda} = (f_i)_{i \in \mathrm{J}_{\lambda}}\) 都是横截的，并以 \(\mathrm{P}_{\lambda}\) 表示该族的纤维积；对于 \(x = (x_i)_{i \in \mathrm{J}_{\lambda}}\) 的每个点 \(\mathrm{P}_{\lambda}\)，S 中的元素 \(f_i(x_i)\) 与 \(i \in \mathrm{J}_{\lambda}\) 无关，记为 \(u_{\lambda}(x)\)；

于是 \(u_{\lambda}\) 是从 \(P_{\lambda}\) 到 S 的态射。族 \(\mathbf{u} = (u_{\lambda})\) 横截的充分必要条件是族 \(\mathbf{f}\) 横截。于是，从 \(\prod_{\lambda \in \Lambda} \prod_{i \in J_{\lambda}} X_i\) 到 \(\prod_{i \in I} X_i\) 的典范双射限制后给出从纤维积 \(\prod_{\lambda \in \Lambda} {}_{S}P_{\lambda}\) 到纤维积 \(\prod_{i \in I} {}_{S}X_i\) 的同构。

5.11.5. 设 S 是一个流形。称 S 上的流形为赋有态射 \(\lambda: X \to S\) 的流形 X。设 \((X, \lambda)\) 是 S 上的流形，\(f: S' \to S\) 是流形的态射，且有序对 \((f, \lambda)\) 横截。记 \(f^*(X)\) 为流形 \(S' \times_{S} X\)，并赋予它由 \(f^*(\lambda): S' \times_{S} X\) 定义的态射，其中 \(f^*(\lambda)(s', x) = s'\)。称 \(f^*(X)\) 是沿 f 通过从 S 到 \(S'\) 的基变换由 X 得到的。若 \(f\)\(\lambda\) 是浸没（分别为浸入、子浸入、étale 态射），则 \(f^*(\lambda)\) 也具有相应性质。

5.11.6. 设 \(f: X \to Y\) 是流形的一个态射，W 是 Y 的子流形；令 \(i\) 为 W 到 Y 的典范嵌入。若有序对 \(f\) 在 \(x \in f^{-1}(W)\) 的点 \((f, i)\) 处横截，则称 \((x, f(x))\) 在点 \(X \times W\) 处横截于 W。其充分必要条件是满足下列条件：

\begin{enumerate}
\def\labelenumi{(\roman{enumi})}
\item
  切空间 \(\mathbf{T}_{f(x)}(\mathbf{Y})\) 是 \(\mathbf{T}_{f(x)}(\mathbf{W})\) 与 \(\mathbf{T}_{x}(f)\) 的像之和；
\item
\(\mathrm{T}_{x}(f)^{-1}(\mathrm{T}_{f(x)}(\mathrm{W}))\) 是 \(f(x)\) 处 W 的切空间的逆像，具有拓扑补空间。
\end{enumerate}

若  在  的每一点处都横截于 ，则称 \(f\) 横截于 \(W\)。\(W\)\(f^{-1}(W)\)

从 X 到 Y 的态射是浸没的充分条件是它横截于 Y 的每一点，而必要条件是它横截于 Y 的每个子流形。有限个态射 \(f_{i}:X_{i}\to S\) 横截的充分必要条件是：由 \(g\) 定义的从 \(\prod_{i\in I}X_{i}\) 到 \(S^{I}\) 的态射 \(g((x_{i})_{i\in I})=(f_{i}(x_{i}))_{i\in I}\) 横截于 \(S^{I}\) 的对角线。

5.11.7. 假设态射 \(f:X\to Y\) 横截于 Y 的子流形 W。则  是 X 的子流形；对于 \(f^{-1}(W)\) 中的 \(x\)，与 \(f^{-1}(W)\) 相切的 \(\mathrm{T}_{x}(\mathrm{X})\) 的子空间是 \(f^{-1}(W)\)\(\mathrm{T}_{x}(f)\) 下 \(\mathrm{T}_{f(x)}(\mathrm{W})\) 的逆像。通过取商，线性映射 \(\mathrm{T}_{x}(f)\) 诱导出从 \(f^{-1}(W)\) 在 \(x\) 处的横截空间到 W 在 \(y=f(x)\) 处的横截空间的拓扑向量空间同构。若 W 在 Y 中 y 处的余维为 \(d\)，则 X 的子流形 \(y\)\(f^{-1}(W)\) 在其每一点处的余维均为 \(d\)。最后，映射 \(f^{-1}(W)\)\(x\mapsto(x,f(x))\) 是从 \(f^{-1}(W)\) 到纤维积 \(\mathrm{X}\times_{\mathrm{Y}}\mathrm{W}\) 的流形同构。

5.11.8. 设 \(Y_{1}\) 和 \(Y_{2}\) 是流形 X 的两个子流形，\(\iota_{j}\) 是 \(Y_{j}\) 到 X 的嵌入。若有序对 \(Y_{1}\) 横截，则称 \(Y_{2}\) 和 \((\iota_{1},\iota_{2})\) 横截；这等价于：对于 \(Y_{1}\cap Y_{2}\) 的每一点 x，\(T_{x}(Y_{1})\) 和 \(T_{x}(Y_{2})\) 作为 \(T_{x}(X)\) 的子空间，其和为 \(T_{x}(X)\)，且它们的交在 \(T_{x}(X)\) 中具有拓扑补空间。在这些条件下，\(Y_{1}\cap Y_{2}\) 是 X 的子流形，并且对于 \(Y_{1}\cap Y_{2}\) 中每个 x，有：

\[
\mathrm{T} _ {x} (\mathrm{Y} _ {1} \cap \mathrm{Y} _ {2}) = \mathrm{T} _ {x} (\mathrm{Y} _ {1}) \cap \mathrm{T} _ {x} (\mathrm{Y} _ {2});
\]

若进一步 \(\mathbf{Y}_i\) 在 \(d_i\) 处的余维为 \(x\)，则 \(\mathbf{Y}_1 \cap \mathbf{Y}_2\) 在 \(d_1 + d_2\) 处的余维为 \(x\)。

5.11.9. 设 \(f\) 和 \(g\) 是从流形 X 到流形 Y 的两个态射。若态射 \((f,g):\mathrm{X}\to\mathrm{Y}\times\mathrm{Y}\) 横截于 \(\mathrm{Y}\times\mathrm{Y}\) 的对角线，则 X 中由满足 \(x\) 的点 \(f(x)=g(x)\) 组成的子集 N 是 X 的子流形；称其为双箭头的核。

\[
f, g: \mathrm{X} \longrightarrow \mathrm{Y}.
\]

